\subsection{I-23. A binary-code cost for preserving decomposable exterior states}
\label{I-23}

\paragraph{The specified next construction.}
One proposed repair is to keep the state of an exterior-power lift
on its Grassmann cone by requiring each Hoza seed to act diagonally
on the Pl\"ucker coordinates in a decomposability-preserving way.
We analyze exactly this requirement.  The statements concern
formal algebraic one-parameter actions over a field, not just
their values on the finitely many elements of a finite field.

Fix $2\le r\le w-2$, index Pl\"ucker coordinates by the
$r$-subsets $I\subseteq[w]$, and let $\rho_I\in\mathbb Z$.
The diagonal action $p_I\mapsto t^{\rho_I}p_I$ preserves all
decomposable $r$-vectors if and only if
\[
 \rho_I=c+\sum_{i\in I}a_i
 \quad\hbox{for some }c,a_1,\ldots,a_w\in\mathbb Z.
 \tag{I.28}\label{eq:I-plucker-additive}
\]
Consequently, if a binary code
$C:\binom{[w]}r\to\{0,1\}^{\ell}$ has minimum distance $d$
and every seed of the Hoza weighting
\[
 \rho_s(I)=\sum_{j=1}^{\ell}6^{\ell-j}C(I)_js_j,
 \qquad s\in\{0,1\}^{\ell},
\]
preserves decomposability, then
\[
 \ell\ge\frac{wd}{2}.
 \tag{I.29}\label{eq:I-plucker-code-cost}
\]
The additional layer multiplier $6^{\nu_2(i)}$ in I-B01
does not change this conclusion.

\paragraph{Proof of the cocharacter classification.}
Let $K$ have size $r-2$, and let $a,b,c,d$ be four different
indices outside $K$.  The Pl\"ucker relation on those indices
forces
\[
 \rho_{Kab}+\rho_{Kcd}
 =\rho_{Kac}+\rho_{Kbd}
 =\rho_{Kad}+\rho_{Kbc}.
\]
Here subscripts denote sets, with the usual orientation signs
in the coordinate relation.  To verify the force of this
condition directly, restrict to planes containing the fixed
coordinate $(r-2)$-plane on $K$.  In the remaining two rows,
use four columns $(e_1,e_2,e_1,e_2)$, then
$(e_1,e_1,e_2,e_2)$.  The relation has exactly two nonzero
terms in each case.  After the diagonal action their two powers
of the indeterminate $t$ must agree.  This gives the displayed
equalities, in every characteristic.

For distinct $i,j$, consider
$\rho_{S\cup\{i\}}-\rho_{S\cup\{j\}}$ with
$|S|=r-1$ and $S\cap\{i,j\}=\varnothing$.
The four-index equality shows that this difference is unchanged
when one element of $S$ is exchanged.  The graph of such subsets
under single exchanges is connected, so the difference is a
constant $d_{ij}$ independent of $S$.  A common $S$ disjoint
from three specified indices exists because $r\le w-2$;
it gives $d_{ij}+d_{jk}=d_{ik}$.  Take $a_1=0$ and
$a_i=d_{i1}$.  Then
$\rho_I-\sum_{i\in I}a_i$ is unchanged by every single exchange
of an $r$-subset.  All $r$-subsets are connected by such
exchanges, so this is one integer $c$.  This proves necessity
of (I.28).  Conversely (I.28) is the exterior action induced by
$\operatorname{diag}(t^{a_1},\ldots,t^{a_w})$, followed by
one common scalar, and plainly preserves decomposable vectors.

\paragraph{Binary-valued additive functions are individual membership bits.}
Suppose $b(I)\in\{0,1\}$ has the form (I.28).
Every exchange difference $a_i-a_j$ lies in $\{-1,0,1\}$.
If $b$ is nonconstant, the $a_i$ therefore have exactly two
levels differing by one.  Let $S$ index the higher level.
Apart from an additive constant, $b(I)=|I\cap S|$.
If both $|S|$ and $w-|S|$ are at least two, then among
$r$-subsets this intersection count ranges from
$\max(0,r-(w-|S|))$ to $\min(r,|S|)$, a range of length
at least two when $2\le r\le w-2$.  That contradicts binary
values.  Hence $|S|=1$ or $w-|S|=1$.
The binary values then force $b(I)$ to be either membership
of one element or its complement.  Constant bits are the only
remaining possibilities.

Apply this classification to the seed $s=e_j$.  The positive
integer factor $6^{\ell-j}$ does not affect the required
exchange equalities, so each code bit is a constant or a single
membership/complement test.  Let $m_i$ count the nonconstant
bits testing element $i$.  Bases differing only by exchange of
$i$ and $j$ have Hamming distance $m_i+m_j$.  Such bases exist
for every pair.  Thus $m_i+m_j\ge d$ for every $i\ne j$.
Summing over all pairs gives
\[
 (w-1)\sum_i m_i\ge\binom w2d,
 \qquad \ell\ge\sum_i m_i\ge wd/2.
\]
This proves (I.29). \hfill$\square$

\paragraph{Consequence for this repair.}
On the $W=\binom wr$ exterior states, using the stated Hoza
error bound $2^{-d}W^4N^2\le1/2$ requires
$d\ge\lceil\log_2(2W^4N^2)\rceil$.
The all-seeds decomposability requirement then forces
$\ell=\Omega(w\log(WN))$ and at least
$(2W^4N^2)^{w/2}$ seeds.  Already $r=2$ gives
$W=\Theta(w^2)$, so this loses the polynomial family size
without invoking an exponentially large exterior dimension.

The restriction is essential: this is not a lower bound for
all generators or all nonlinear decodings of Hoza seeds.
It does not rule out a separately proved successful subset
of coherent seeds from another code, or a construction that
leaves the Grassmann cone at intermediate stages and restores
the required coefficient information afterwards.  Nor would
mere Pl\"ucker compatibility by itself resolve the distinct
variable-quotient condition in I-02.
